% =====================================================================
%  main.tex --- START HERE
%
%  A working deck, already set up. Change the title, replace the
%  slides, and you have your talk. Everything below is meant to be
%  edited or deleted.
%
%  Build:  pdflatex main    (twice -- the outline slides need a second
%                            pass; a third if you add citations)
%
%  demo.pdf is the manual: every component with the code beside it.
%  Snippets to paste live in snippets/. Setup notes are in README.md.
% =====================================================================

\documentclass[11pt,aspectratio=43]{beamer}   % 169 for widescreen
\usetheme{lucid}

% Useful switches -- uncomment one at a time to see what it does:
%   \usetheme[palette=maroon]{lucid}   different accent colour
%   \usetheme[notes=pages]{lucid}      print your speaker notes
%   \usetheme[hidden=show]{lucid}      show the parked slides
%   \usetheme[blind]{lucid}            hide author and institution
% and add "handout" to \documentclass to flatten animations for printing.

\usepackage{tikz}       % only if you draw things
% \usepackage{listings} % only if you show code

\title{Your Title Here}
\subtitle{An optional second line}
\lucidevent{Venue, date}
\date{}

% Nothing identifying appears unless you set it. Uncomment to switch on:
% \author{Your Name}
% \institute{Your Institution}
% \lucidlogo{logo.pdf}

\begin{document}

% --------------------------------------------------------------------
\begin{frame}[plain]
  \titlepage
\end{frame}

% --------------------------------------------------------------------
\section{The problem}

\begin{frame}{What is broken}
  \begin{itemize}
    \item State the problem in one line
    \item Say who has it and why it matters
  \end{itemize}
  \citeline{A small source line, if you need one.}
\end{frame}

\begin{frame}{Why it is hard}
  \begin{steps}
    \item The obvious approach fails because \ldots
    \item The next one fails because \ldots
    \item Which leaves \term{the gap you are filling}
  \end{steps}
\end{frame}

% --------------------------------------------------------------------
\section{What we did}

\begin{frame}{The idea}
  \begin{figuretext}[0.55]
    \begin{itemize}
      \item One sentence of method
      \item One sentence of why it works
    \end{itemize}
    \nextpane
    % \slidefigure{diagram.pdf}
    \begin{tikzpicture}[scale=0.5]
      \draw[lucidMuted] (0,0) rectangle (5,3.6);
      \node[lucidMuted,font=\tiny] at (2.5,1.8) {your figure here};
    \end{tikzpicture}
  \end{figuretext}
\end{frame}

\begin{frame}{The one equation}
  \begin{itemize}
    \item Say what it means before you show it
      \[ \mathcal{L} = \mathbb{E}\left[ \sum_t \gamma^t r_t \right] \]
  \end{itemize}
\end{frame}

% --------------------------------------------------------------------
\section{What happened}

\begin{frame}{Results}
  \begin{lucidtable}
    \begin{tabular}{lrr}
      \toprule
      \thead{Method & Setting A & Setting B} \\
      \midrule
      Baseline & 62 & 55 \\
      Ours     & \term{78} & \term{71} \\
      \bottomrule
    \end{tabular}
  \end{lucidtable}
  \speakernote{Say the numbers out loud. If asked about variance,
    go to the backup slide.}
\end{frame}

\begin{frame}{Takeaway}
  \takeaway{The one sentence you want them to leave with.}
\end{frame}

% --------------------------------------------------------------------
% Parked slides. Nothing in here is compiled unless you build with
% \usetheme[hidden=show]{lucid} -- so it is a safe place to leave things.
\begin{hidden}
\begin{frame}{Cut for time}
  \begin{itemize}
    \item Still here, just not in the deck
  \end{itemize}
\end{frame}
\end{hidden}

% --------------------------------------------------------------------
% Backup slides. These do not count toward the page number.
\backupframes

\begin{frame}{Backup: extra detail}
  \begin{itemize}
    \item The thing someone always asks about
  \end{itemize}
\end{frame}

\end{document}
